\documentclass[UTF8,fontset=windows]{ctexart}
\usepackage[a4paper,margin=22mm]{geometry}
\usepackage{amsmath,amssymb}
\usepackage{tikz}
\usepackage{booktabs}
\usepackage{enumitem}
\usepackage{xcolor}
\usepackage{float}
\usetikzlibrary{arrows.meta,calc,patterns}

\setlength{\parindent}{0pt}
\setlength{\parskip}{6pt}
\setlist[itemize]{leftmargin=2em,itemsep=2pt,topsep=2pt}

\definecolor{accent}{RGB}{32,92,148}
\definecolor{softblue}{RGB}{224,238,250}
\definecolor{softred}{RGB}{252,230,226}
\definecolor{axisgray}{RGB}{80,80,80}

\newcommand{\diff}{\,\mathrm{d}}

\title{\vspace{-1.5cm}\textbf{清华大学 2024 年《材料力学》（吴坚班）期末考试解答}}
\author{回忆整理版标准解答}
\date{}

\begin{document}
\maketitle
\vspace{-1.0cm}

\section*{一、微元强度与应变基础题}

\textbf{1、微元体处于双轴拉伸状态，拉应力分别为 $\sigma_x$ 和 $\sigma_y$，材料许用应力为 $[\sigma]$，进行强度校核。}

\subsubsection*{解答：}
主应力为 $\sigma_1 = \sigma_x$，$\sigma_2 = \sigma_y$，$\sigma_3 = 0$（假设双轴均拉，且 $\sigma_x > \sigma_y > 0$）。
根据第三强度理论（Tresca 屈服准则）：
\[
    \sigma_{r3} = \sigma_1 - \sigma_3 = \sigma_x \le [\sigma]
\]

\textbf{2、微元发生剪切变形，已知变形前直角为 $\alpha = \pi/2$，变形后该角变为 $\beta$，求剪应变 $\gamma$。}

\subsubsection*{解答：}
剪应变定义为直角的改变量：
\[
    \gamma = \frac{\pi}{2} - \beta
\]

\textbf{3、宽度为 $b$，厚度为 $h$ 的平板条，无管弯曲卷入半径为 $R$ 的圆筒中，已知材料弹性模量为 $E$。求最大弯曲正应力。}

\subsubsection*{解答：}
弯曲曲率为 $\kappa = 1/R$，由弯曲正应力公式：
\[
    \sigma = \frac{E \cdot y}{R} \implies \sigma_{\max} = \frac{E h}{2R}
\]

\newpage

\section*{二、悬臂圆管弯扭组合结构题}

\textbf{4、水平悬臂圆管 AB 长度为 $L$，A 端固定。在自由端 B 处连接一水平拐角折杆 BC，长度为 $L$（平行于 $y$ 轴）。C 端受轴向拉力 $F_x$ 和垂直向下的力 $F_z$。\\
（1）画出 AB 段的弯矩图和扭矩图；\\
（2）确定危险截面和危险点；\\
（3）按第三强度理论对圆管进行强度校核。}

\subsubsection*{解答：}
\textbf{（1）内力图分析}
\begin{itemize}
    \item 扭矩：$F_z$ 在 C 端对 AB 轴线产生的扭矩恒为：$T = F_z \cdot L$。
    \item 弯矩：
    $F_z$ 引起绕 $y$ 轴的弯矩，固定端最大：$M_y(x) = F_z(L-x)$。
    $F_x$ 沿 $x$ 轴，偏心矩为 $L$（沿 $y$ 方向），引起绕 $z$ 轴的弯矩：$M_z = F_x \cdot L$。
\end{itemize}
因此，固定端 A 处的合成弯矩为：
\[
    M_{\text{sum}} = \sqrt{(F_z L)^2 + (F_x L)^2} = L\sqrt{F_x^2 + F_z^2}
\]
扭矩为 $T = F_z L$。

\textbf{（2）危险截面和危险点}
\begin{itemize}
    \item 危险截面：固定端 A 处 ($x=0$)。
    \item 危险点：外表面弯曲拉应力最大处（由弯曲矢量确定）。
\end{itemize}

\textbf{（3）强度校核}
当量弯矩为：
\[
    M_{r3} = \sqrt{M_{\text{sum}}^2 + T^2} = L\sqrt{F_x^2 + 2F_z^2}
\]
强度条件：
\[
    \sigma_{r3} = \frac{32 M_{r3}}{\pi d^3 (1 - \alpha^4)} \le [\sigma] \quad (\alpha = d_{\text{in}}/d_{\text{out}})
\]

\newpage

\section*{三、组合梁内力与位移计算题}

\textbf{5、简支外伸梁如图所示，AB 段长为 $2L$，承受均布载荷 $q$，BC 段长为 $L$，C 端承受集中力 $q L$。\\
（1）计算 A、B 支座的反力；\\
（2）绘制剪力图和弯矩图；\\
（3）确定危险截面；\\
（4）计算 B 点转角与 C 点挠度。}

\subsubsection*{解答：}
\textbf{（1）反力计算}
设 $R_A, R_B$ 向上。
\[
    \sum M_A = 0 \implies R_B \cdot 2L - q(2L) \cdot L - qL \cdot 3L = 0 \implies R_B = 2.5 qL
\]
\[
    R_A = 3qL - R_B = 0.5 qL
\]

\textbf{（2）内力图}
剪力在 $x = 0.5L$ 处为零，该处弯矩最大 $M = \frac{1}{8}qL^2$。B点截面弯矩为 $M_B = -qL^2$。

\textbf{（3）危险截面}
弯矩最大处在 \textbf{支座 B 截面}。

\textbf{（4）位移计算}
使用积分法或叠加法求得：
\[
    \theta_B = -\frac{qL^3}{3EI}, \quad w_C = \frac{17 q L^4}{12 EI}
\]

\newpage

\section*{四、薄壁圆筒受扭受压强度题}

\textbf{6、薄壁圆筒直径为 $D$，壁厚为 $h$，内腔受压力 $P$ 作用，两端受扭矩 $M_x$ 作用。\\
（1）校核圆筒强度；\\
（2）求表面任意点与轴线呈 $60^\circ$ 方向的应力。}

\subsubsection*{解答：}
\textbf{（1）圆筒主应力计算}
\begin{itemize}
    \item 环向正应力：$\sigma_t = \frac{P D}{2h}$。
    \item 轴向正应力：$\sigma_l = \frac{P D}{4h}$。
    \item 扭转剪应力：$\tau = \frac{2 M_x}{\pi D^2 h}$。
\end{itemize}
主应力为：
\[
    \sigma_{1,2} = \frac{\sigma_t + \sigma_l}{2} \pm \sqrt{\left(\frac{\sigma_t - \sigma_l}{2}\right)^2 + \tau^2}
\]
代入第三或第四强度理论公式即可。

\textbf{（2）$60^\circ$ 方向的应力}
利用坐标变换公式：
\[
    \sigma_{60^\circ} = \sigma_l \cos^2(60^\circ) + \sigma_t \sin^2(60^\circ) + 2\tau \sin(60^\circ)\cos(60^\circ)
\]

\newpage

\section*{五、并联受力与失稳分析题}

\textbf{7、刚性梁 AB 长度为 $2L$，A 端为固定铰支座。在梁的 C 点（距 A 为 $L$）和 B 点（距 A 为 $2L$）分别连接两根倾斜为 $45^\circ$ 的钢拉杆 CD 和 BE，两杆的直径分别为 $d_1, d_2$。在 B 端施加竖向向下集中荷载 $F$。已知材料弹性模量 $E$，屈服极限 $\sigma_s$，安全系数 $n_{st}$，求许用载荷 $F$。}

\subsubsection*{解答：}
由于梁为刚性杆，其转动位移决定了两拉杆的变形关系。
\begin{itemize}
    \item 变形协调：C 点的垂直位移是 B 点的一半。
    \item 轴力分配：设拉杆轴力为 $N_1, N_2$，则：
    \[
        N_1 = \frac{E A_1 \Delta_C}{L_1}, \quad N_2 = \frac{E A_2 \Delta_B}{L_2} \implies N_2 = 2 N_1 \frac{A_2}{A_1}
    \]
    \item 矩平衡：由对 A 的力矩平衡解出 $F$。
    \item 考虑屈服与失稳安全限制条件确定 $F$。
\end{itemize}

\end{document}
